\documentclass[11pt]{article}
\usepackage[margin=1in]{geometry}
\usepackage{amsmath,amssymb}
\usepackage{enumitem}

\title{MATH 407: Homework 2}
\author{Yazan Ismail}
\date{October 18, 2026}

\begin{document}
\maketitle

%==================================================
\section*{Exercise 1}

\noindent\textbf{Problem.}
Find all vertices of the polyhedron of points $(x,y,z) \in \mathbb{R}^3$ satisfying the following 5 constraints:
\[
\text{(1) } x \ge 0 \qquad
\text{(2) } y \ge 0 \qquad
\text{(3) } z \ge 0 \qquad
\text{(4) } x + 2y \le 4 \qquad
\text{(5) } x + 2y + 3z \le 6
\]
For each vertex $\mathbf{v}$, write down the subset of $\{1,2,3,4,5\}$ indexing the set of inequalities for which equality holds at $\mathbf{v}$.

For example, for the polyhedron of points $(x,y) \in \mathbb{R}^2$ satisfying
\[
\text{(1) } x \ge 0 \qquad
\text{(2) } y \ge 0 \qquad
\text{(3) } x + y \le 1,
\]
the analogous answer would be:
\begin{itemize}
  \item $(0,0) \leftrightarrow \{1,2\}$ -- I.e.\ there is a vertex at $(0,0)$ and constraints 1 and 2 hold with an equality there.
  \item $(0,1) \leftrightarrow \{1,3\}$ -- I.e.\ there is a vertex at $(0,1)$ and constraints 1 and 3 hold with an equality there.
  \item $(1,0) \leftrightarrow \{2,3\}$ -- I.e.\ there is a vertex at $(1,0)$ and constraints 2 and 3 hold with an equality there.
\end{itemize}
It is not necessary to detail how you solve each system of equations.

\medskip
\noindent\textbf{Solution.}
For each of the $\binom{5}{3} = 10$ choices of three constraints, I set those three constraints to equality, checked that their coefficient rows are linearly independent, solved the resulting system, and checked the solution against all five constraints. The vertices are:
\begin{align*}
(0,0,0) &\leftrightarrow \{1,2,3\} \\
(0,0,2) &\leftrightarrow \{1,2,5\} \\
(0,2,0) &\leftrightarrow \{1,3,4\} \\
(0,2,\tfrac{2}{3}) &\leftrightarrow \{1,4,5\} \\
(4,0,0) &\leftrightarrow \{2,3,4\} \\
(4,0,\tfrac{2}{3}) &\leftrightarrow \{2,4,5\}
\end{align*}
At each of these points, no constraint outside the listed set holds with equality, so each listed set is the complete set of tight constraints.

The remaining four choices do not give vertices:
\begin{itemize}
  \item $\{1,2,4\}$: dependent. Setting $x = y = 0$ forces $x + 2y = 0 \neq 4$, so the system has no solution.
  \item $\{3,4,5\}$: dependent. With $z = 0$, constraints (4) and (5) become $x + 2y = 4$ and $x + 2y = 6$, which is inconsistent.
  \item $\{1,3,5\}$: gives $(0,3,0)$, which violates (4) since $0 + 2(3) = 6 > 4$.
  \item $\{2,3,5\}$: gives $(6,0,0)$, which violates (4) since $6 > 4$.
\end{itemize}

%==================================================
\clearpage
\section*{Exercise 2}

\noindent\textbf{Problem.}
Consider the linear program
\begin{align*}
\max\quad & x_1 - x_2 + x_4 - x_5 \\
\text{subject to}\quad & x_1 + x_2 + x_3 \le 5 \\
& x_3 + x_5 \le 3 \\
& x_4 + x_5 \le 3 \\
& x_1, x_2, x_3, x_4, x_5 \ge 0
\end{align*}
and the point $\mathbf{v} = (3, 2, 0, 0, 3)$.
\begin{enumerate}[label=(\alph*)]
  \item Find a nonzero vector $\mathbf{w}$ so that $\mathbf{v} + \mathbf{w}$ and $\mathbf{v} - \mathbf{w}$ both belong to the feasible region of this LP. (This shows that $\mathbf{v}$ is not a vertex!)
  \item For what values of $\lambda \in \mathbb{R}$ is $\mathbf{v} + \lambda \mathbf{w}$ feasible, where $\mathbf{w}$ is as in the previous item?
  \item Which of the feasible points $\mathbf{v} + \lambda \mathbf{w}$ (that you found in part (b)) achieves the largest value of the objective function?
\end{enumerate}

\medskip
\noindent\textbf{Solution.}
\begin{enumerate}[label=(\alph*)]
  \item $\mathbf{w} = (1, -1, 0, 0, 0)$.

  At $\mathbf{v}$, the constraints that hold with equality are $x_1 + x_2 + x_3 \le 5$, $x_3 + x_5 \le 3$, $x_4 + x_5 \le 3$, $x_3 \ge 0$, and $x_4 \ge 0$. For both $\mathbf{v} + \mathbf{w}$ and $\mathbf{v} - \mathbf{w}$ to be feasible, $\mathbf{w}$ must leave each of these tight constraints unchanged:
  \[
  w_1 + w_2 + w_3 = 0, \qquad
  w_3 + w_5 = 0, \qquad
  w_4 + w_5 = 0, \qquad
  w_3 = 0, \qquad
  w_4 = 0.
  \]
  This gives $w_3 = w_4 = w_5 = 0$ and $w_2 = -w_1$. Taking $w_1 = 1$ gives $\mathbf{v} + \mathbf{w} = (4, 1, 0, 0, 3)$ and $\mathbf{v} - \mathbf{w} = (2, 3, 0, 0, 3)$, and both are feasible.

  \item $\lambda \in [-3, 2]$.

  We have $\mathbf{v} + \lambda \mathbf{w} = (3 + \lambda,\ 2 - \lambda,\ 0,\ 0,\ 3)$. The three main constraints hold for every $\lambda$:
  \[
  (3 + \lambda) + (2 - \lambda) + 0 = 5 \le 5, \qquad
  0 + 3 \le 3, \qquad
  0 + 3 \le 3.
  \]
  Nonnegativity requires $3 + \lambda \ge 0$ and $2 - \lambda \ge 0$, so $-3 \le \lambda \le 2$.

  \item $\lambda = 2$, giving the point $\mathbf{v} + 2\mathbf{w} = (5, 0, 0, 0, 3)$ with objective value $2$.

  Along this line, the objective is
  \[
  (3 + \lambda) - (2 - \lambda) + 0 - 3 = 2\lambda - 2,
  \]
  which is increasing in $\lambda$. It is therefore largest at the upper end of the range from part (b), $\lambda = 2$, where its value is $2(2) - 2 = 2$.
\end{enumerate}

%==================================================
\clearpage
\section*{Exercise 3}

\noindent\textbf{Problem.}
Reformulate the following linear program in equational form.
\begin{align*}
\min\quad & x_1 + x_2 - x_3 + x_4 - x_5 + x_6 \\
\text{subject to}\quad & x_1 + x_5 - 3x_6 \ge 5 \\
& x_3 + x_5 \le 3 \\
& x_2 + x_3 + x_5 \ge 6 \\
& x_2 + x_4 + x_6 \le 3 \\
& x_1, x_2, x_3, x_4, x_5 \ge 0
\end{align*}

\medskip
\noindent\textbf{Solution.}
I converted the minimization into a maximization by negating the objective, and multiplied the two $\ge$ constraints by $-1$. Since $x_6$ is not constrained in sign, I substituted $x_6 = y_6 - z_6$ with $y_6, z_6 \ge 0$. Finally, I added a nonnegative slack variable $s_i$ to each constraint. The equational form is
\begin{align*}
\max\quad & -x_1 - x_2 + x_3 - x_4 + x_5 - (y_6 - z_6) \\
\text{subject to}\quad & -x_1 - x_5 + 3(y_6 - z_6) + s_1 = -5 \\
& x_3 + x_5 + s_2 = 3 \\
& -x_2 - x_3 - x_5 + s_3 = -6 \\
& x_2 + x_4 + (y_6 - z_6) + s_4 = 3 \\
& x_1, x_2, x_3, x_4, x_5, y_6, z_6, s_1, s_2, s_3, s_4 \ge 0
\end{align*}

\end{document}
